\documentclass[11pt]{article}

\usepackage[margin=1in]{geometry}
\usepackage[T1]{fontenc}
\usepackage[utf8]{inputenc}
\usepackage{amsmath}
\usepackage{array}
\usepackage{booktabs}
\usepackage{graphicx}
\usepackage{hyperref}
\usepackage{listings}
\usepackage{longtable}
\usepackage{xcolor}

\hypersetup{
  colorlinks=true,
  linkcolor=blue,
  urlcolor=blue,
  citecolor=blue
}

\lstset{
  basicstyle=\ttfamily\small,
  breaklines=true,
  frame=single,
  columns=fullflexible
}

\title{RES201 Week 2 Report\\Multi-Seed Fine-Tuning Learning Curves for Oxides and Nitrides}
\author{Prepared from workspace execution in \texttt{res201-alignn-domain-shift}}
\date{April 8, 2026}

\begin{document}
\maketitle

\section{Abstract}
This report documents the complete Week 2 execution for the RES201 project
\textit{Chemical-Family Domain Shift in Pretrained ALIGNN Data Efficiency for Oxide vs. Nitride Formation Energy Prediction}.
The Week 2 goal was to complete the fine-tuning learning-curve study for oxide and nitride subsets using the pretrained ALIGNN formation-energy model, then report the corresponding error tables and plots.
Following the broader project requirements, the final workspace state now includes three random seeds for each family and each training size.

The completed Week 2 study therefore contains 36 fine-tuning runs in total:
two chemical families, six training sizes, and three seeds.
All aggregate tables and learning-curve plots have been generated.
The main scientific outcome is that oxide adaptation is smoother and less variable than nitride adaptation, but under the current partial-fine-tuning protocol neither family outperforms its zero-shot baseline by $N=1000$.

\section{Week 2 Objective and Deliverables}
According to the RES201 project brief, Week 2 required:
\begin{itemize}
  \item all oxide fine-tuning runs completed;
  \item an oxide MAE table;
  \item an oxide learning-curve plot;
  \item all nitride fine-tuning runs completed;
  \item a nitride MAE table; and
  \item a nitride learning-curve plot.
\end{itemize}

The same brief also states that the overall project should use multiple random seeds, with three as a minimum target and five as a preferred target.
For that reason, this Week 2 report documents the final three-seed study rather than only a single seed.

\section{Project Context}
\subsection{Scientific Question}
The project investigates whether a pretrained ALIGNN model transfers differently to oxide and nitride subsets of the JARVIS \texttt{dft\_3d\_2021} formation-energy dataset.
The motivating hypothesis is that nitride materials are more out-of-domain relative to the chemistry emphasized during pretraining, and therefore may require more in-family data or more flexible adaptation to achieve the same performance as oxides.

\subsection{Family Definitions}
The family definitions inherited from the Week 1 stage-2 dataset pipeline are:
\begin{itemize}
  \item oxide arm: materials containing oxygen;
  \item nitride arm: materials containing nitrogen and not oxygen;
  \item oxynitrides: excluded from the nitride arm;
  \item original benchmark split IDs applied before family filtering.
\end{itemize}

The resulting frozen family dataset counts used throughout Week 2 are:
\begin{center}
\begin{tabular}{lrrrrr}
\toprule
Family & All & Train & Val & Test & Pool \\
\midrule
Oxide   & 14991 & 11960 & 1547 & 1484 & 13507 \\
Nitride &  2288 &  1837 &  209 &  242 &  2046 \\
\bottomrule
\end{tabular}
\end{center}

These counts match the validated family summaries in the workspace and confirm that the test sets remained fixed while Week 2 varied only the fine-tuning subset size and seed.

\section{Methodology}
\subsection{Pretrained ALIGNN Starting Point}
All Week 2 runs start from the same pretrained formation-energy checkpoint:
\begin{itemize}
  \item \texttt{jv\_formation\_energy\_peratom\_alignn/checkpoint\_300.pt}
  \item \texttt{jv\_formation\_energy\_peratom\_alignn/config.json}
\end{itemize}

The workspace uses ALIGNN as a graph neural network for crystalline materials, with atom-bond structure and angle-sensitive message passing inherited from the pretrained model configuration.
Week 2 does not alter the underlying architecture between oxide and nitride runs; only the fine-tuning subsets differ.

\subsection{Fine-Tuning Protocol}
The fine-tuning strategy used in this workspace is deliberately controlled and conservative:
\begin{enumerate}
  \item initialize the model from the local pretrained checkpoint;
  \item freeze all parameters at the start of fine-tuning;
  \item unfreeze only the final two parameter groups, \texttt{gcn\_layers.3} and \texttt{fc};
  \item fine-tune for 50 epochs with learning rate $10^{-4}$ and batch size 16;
  \item select the best validation checkpoint and evaluate on the fixed test set.
\end{enumerate}

This is the same transfer setting introduced in Week 1, which means the Week 2 comparison isolates the effect of increasing training-set size rather than simultaneously changing multiple training choices.

\subsection{Training Sizes and Seeds}
For each family, the following training sizes were evaluated:
\[
N \in \{10, 50, 100, 200, 500, 1000\}.
\]

For each of those $N$ values, the study uses seeds 0, 1, and 2.
That yields:
\[
2 \text{ families} \times 6 \text{ sizes} \times 3 \text{ seeds} = 36 \text{ runs}.
\]

The per-run train and validation sizes follow the same rule used in Week 1:
\begin{itemize}
  \item $N=10 \Rightarrow 5$ train and $5$ validation;
  \item $N=50 \Rightarrow 45$ train and $5$ validation;
  \item $N=100 \Rightarrow 90$ train and $10$ validation;
  \item $N=200 \Rightarrow 180$ train and $20$ validation;
  \item $N=500 \Rightarrow 450$ train and $50$ validation;
  \item $N=1000 \Rightarrow 900$ train and $100$ validation.
\end{itemize}

\subsection{Dataset Materialization}
Each individual run has its own local dataset root under:
\begin{itemize}
  \item \texttt{results/oxide/N...\_seed.../dataset\_root/}
  \item \texttt{results/nitride/N...\_seed.../dataset\_root/}
\end{itemize}

Each dataset root contains:
\begin{itemize}
  \item \texttt{id\_prop.csv} with the structure IDs and target values;
  \item \texttt{split\_manifest.json} with the exact train, validation, and test membership for that run;
  \item POSCAR files for the selected structures.
\end{itemize}

This design makes every run reproducible and auditable after the fact.

\subsection{Execution Scripts}
The Week 2 workflow in this workspace is built around the following scripts:
\begin{itemize}
  \item \texttt{scripts/shared/run\_week2\_finetune\_suite.py}
  \item \texttt{scripts/shared/summarize\_week2\_finetune.py}
  \item \texttt{scripts/shared/check\_week2\_status.sh}
  \item \texttt{scripts/shared/finetune\_last2\_alignn.py}
  \item \texttt{scripts/shared/prepare\_week1\_finetune\_dataset.py}
\end{itemize}

Together these scripts:
\begin{enumerate}
  \item prepare the correct dataset root for each family, size, and seed;
  \item write the run-specific ALIGNN config JSON;
  \item launch the 50-epoch fine-tuning job;
  \item save histories, prediction CSVs, checkpoints, and summary JSON files;
  \item aggregate the full study into summary CSV tables and learning-curve plots;
  \item verify final completion through a single checker command.
\end{enumerate}

\subsection{Compute Environment}
The Week 2 runs were executed in the local project environment rather than a cloud GPU setup.
In this workspace the fine-tuning commands were run on CPU.
This matters for interpretation because absolute runtime is longer than it would be on a GPU, but it does not affect the logic of the Week 2 deliverables.

The measured runtime metadata recorded in the run summaries total approximately:
\begin{itemize}
  \item 3.65 CPU-hours for oxide runs with recorded timing;
  \item 1.55 CPU-hours for nitride runs with recorded timing;
  \item 5.20 CPU-hours combined.
\end{itemize}

Two early seed-0 $N=50$ summaries were created before elapsed-time logging was added to the fine-tuning wrapper, so those two completed runs do not include runtime fields even though their result artifacts are present.

\section{Execution Procedure}
\subsection{Representative Commands}
The additional seed sweep for Week 2 was executed using commands of the following form:

\begin{lstlisting}[language=bash]
source "$(conda info --base)/etc/profile.d/conda.sh"
conda activate res201_train
python -u scripts/shared/run_week2_finetune_suite.py \
  --repo-root . \
  --families oxide nitride \
  --Ns 10 50 100 200 500 1000 \
  --seeds 1 2 \
  --device cpu \
  --continue-on-error
\end{lstlisting}

Aggregation into CSV summaries and plots was done with:

\begin{lstlisting}[language=bash]
source "$(conda info --base)/etc/profile.d/conda.sh"
conda activate res201_train
python -u scripts/shared/summarize_week2_finetune.py \
  --repo-root . \
  --families oxide nitride \
  --Ns 10 50 100 200 500 1000 \
  --seeds 0 1 2 \
  --out-dir reports/week2
\end{lstlisting}

Final verification is:

\begin{lstlisting}[language=bash]
scripts/shared/check_week2_status.sh .
\end{lstlisting}

The run-suite summary file reports that the seed-1 and seed-2 completion batch requested 24 runs, completed 24 runs, and had 0 failures.

\section{Results}
\subsection{Zero-Shot Baselines}
The fixed zero-shot reference MAEs used for comparison are:

\begin{center}
\begin{tabular}{lrr}
\toprule
Family & Test set size & Zero-shot MAE (eV/atom) \\
\midrule
Oxide   & 1484 & 0.03418 \\
Nitride &  242 & 0.06954 \\
\bottomrule
\end{tabular}
\end{center}

\subsection{Aggregate Multi-Seed Results}
The summary table below reports the mean and standard deviation of test MAE across seeds 0, 1, and 2 for each family and each training size.

\begin{center}
\small
\begin{tabular}{lrrrrrrrrr}
\toprule
Family & $N$ & Runs & Train & Val & Test & Mean MAE & Std & Zero-shot & Gain \\
\midrule
Nitride & 10   & 3 &   5 &   5 &  242 & 0.09875 & 0.01763 & 0.06954 & -0.02921 \\
Nitride & 50   & 3 &  45 &   5 &  242 & 0.14816 & 0.02210 & 0.06954 & -0.07862 \\
Nitride & 100  & 3 &  90 &  10 &  242 & 0.23754 & 0.06205 & 0.06954 & -0.16800 \\
Nitride & 200  & 3 & 180 &  20 &  242 & 0.17524 & 0.06549 & 0.06954 & -0.10570 \\
Nitride & 500  & 3 & 450 &  50 &  242 & 0.11054 & 0.00372 & 0.06954 & -0.04100 \\
Nitride & 1000 & 3 & 900 & 100 &  242 & 0.10040 & 0.00262 & 0.06954 & -0.03086 \\
Oxide   & 10   & 3 &   5 &   5 & 1484 & 0.04618 & 0.01313 & 0.03418 & -0.01200 \\
Oxide   & 50   & 3 &  45 &   5 & 1484 & 0.06174 & 0.00997 & 0.03418 & -0.02756 \\
Oxide   & 100  & 3 &  90 &  10 & 1484 & 0.05272 & 0.00166 & 0.03418 & -0.01854 \\
Oxide   & 200  & 3 & 180 &  20 & 1484 & 0.05174 & 0.00337 & 0.03418 & -0.01756 \\
Oxide   & 500  & 3 & 450 &  50 & 1484 & 0.04748 & 0.00107 & 0.03418 & -0.01330 \\
Oxide   & 1000 & 3 & 900 & 100 & 1484 & 0.04556 & 0.00086 & 0.03418 & -0.01138 \\
\bottomrule
\end{tabular}
\normalsize
\end{center}

The final column reports transfer gain relative to zero-shot, defined here as:
\[
\text{gain} = \text{zero-shot MAE} - \text{mean fine-tuned test MAE}.
\]
Every gain value is negative in Week 2, which means that under the current protocol the fine-tuned mean remains worse than zero-shot for every family and every $N$.

\subsection{Best Observed Outcomes}
The most favorable outcomes observed in Week 2 are:
\begin{itemize}
  \item oxide best mean fine-tuned MAE: $0.04556 \pm 0.00086$ at $N=1000$;
  \item nitride best mean fine-tuned MAE: $0.09875 \pm 0.01763$ at $N=10$;
  \item oxide best single run: 0.03604 at $N=10$, seed 2;
  \item nitride best single run: 0.08263 at $N=10$, seed 1.
\end{itemize}

Even the best single oxide run is still slightly worse than the oxide zero-shot baseline of 0.03418, and the best nitride run is still worse than the nitride zero-shot baseline of 0.06954.

\subsection{Interpretation of the Learning Curves}
Several important patterns emerge from the Week 2 results:
\begin{enumerate}
  \item \textbf{Oxide improves steadily with more data.} Although oxide does not beat zero-shot, its mean MAE drops from 0.06174 at $N=50$ to 0.04556 at $N=1000$.
  \item \textbf{Oxide variance shrinks strongly at larger $N$.} The oxide standard deviation falls to 0.00086 at $N=1000$, indicating a stable response to additional in-family data.
  \item \textbf{Nitride is substantially noisier.} Nitride shows large seed-to-seed variation at intermediate data sizes, especially at $N=100$ and $N=200$.
  \item \textbf{Nitride remains harder to adapt.} Even when nitride variance becomes smaller at $N=500$ and $N=1000$, the absolute MAE still remains far above the nitride zero-shot baseline.
  \item \textbf{The domain-shift hypothesis remains plausible.} The oxide arm responds more smoothly to increasing fine-tuning data, whereas the nitride arm behaves more erratically, which is consistent with stronger chemistry mismatch relative to pretraining.
  \item \textbf{The current transfer setting may be too restrictive.} Because only the final two parameter groups are trainable, the model may not have enough flexibility to adapt effectively to the family-restricted distributions.
\end{enumerate}

One noteworthy nitride result is that the best mean appears at $N=10$ rather than $N=1000$.
This should not be interpreted as evidence that less nitride data is inherently preferable.
It is more likely a sign that the present transfer configuration is unstable or under-parameterized for this family.

\section{Plots}
The final Week 2 report should include the oxide and nitride learning curves.
This LaTeX file is intentionally written so it still compiles even if the plot files are not present in the Overleaf project yet.
If the files are present, they are inserted below.

\begin{figure}[htbp]
\centering
\IfFileExists{reports/week2/oxide_learning_curve.pdf}{
  \includegraphics[width=0.82\linewidth]{reports/week2/oxide_learning_curve.pdf}
}{
  \IfFileExists{week2/oxide_learning_curve.pdf}{
    \includegraphics[width=0.82\linewidth]{week2/oxide_learning_curve.pdf}
  }{
    \fbox{\parbox{0.82\linewidth}{\centering
      Oxide learning-curve figure placeholder.\\
      Upload \texttt{reports/week2/oxide\_learning\_curve.pdf} or\\
      \texttt{week2/oxide\_learning\_curve.pdf} to display the plot here.}}
  }
}
\caption{Oxide learning curve across $N=\{10,50,100,200,500,1000\}$ using three seeds.}
\end{figure}

\begin{figure}[htbp]
\centering
\IfFileExists{reports/week2/nitride_learning_curve.pdf}{
  \includegraphics[width=0.82\linewidth]{reports/week2/nitride_learning_curve.pdf}
}{
  \IfFileExists{week2/nitride_learning_curve.pdf}{
    \includegraphics[width=0.82\linewidth]{week2/nitride_learning_curve.pdf}
  }{
    \fbox{\parbox{0.82\linewidth}{\centering
      Nitride learning-curve figure placeholder.\\
      Upload \texttt{reports/week2/nitride\_learning\_curve.pdf} or\\
      \texttt{week2/nitride\_learning\_curve.pdf} to display the plot here.}}
  }
}
\caption{Nitride learning curve across $N=\{10,50,100,200,500,1000\}$ using three seeds.}
\end{figure}

\section{Week 2 Deliverable Checklist}
\begin{longtable}{p{0.43\linewidth}p{0.10\linewidth}p{0.35\linewidth}}
\toprule
Deliverable & Status & Evidence \\
\midrule
All oxide fine-tuning runs completed & Done & \texttt{results/oxide/N10\_seed0} through \texttt{results/oxide/N1000\_seed2} \\
All nitride fine-tuning runs completed & Done & \texttt{results/nitride/N10\_seed0} through \texttt{results/nitride/N1000\_seed2} \\
Oxide MAE table & Done & embedded in this report and saved in \texttt{reports/week2/finetune\_summary\_by\_N.csv} \\
Nitride MAE table & Done & embedded in this report and saved in \texttt{reports/week2/finetune\_summary\_by\_N.csv} \\
Oxide learning-curve plot & Done & \texttt{reports/week2/oxide\_learning\_curve.pdf} and \texttt{.png} \\
Nitride learning-curve plot & Done & \texttt{reports/week2/nitride\_learning\_curve.pdf} and \texttt{.png} \\
Three-seed reporting & Done & every family and $N$ entry in the aggregate CSV has \texttt{runs = 3} \\
\bottomrule
\end{longtable}

\section{Artifact Map}
The most important Week 2 files and folders in the workspace are:
\begin{itemize}
  \item \texttt{configs/oxide\_week2\_N...\_seed....finetune\_last2.json}: oxide run-specific ALIGNN configs.
  \item \texttt{configs/nitride\_week2\_N...\_seed....finetune\_last2.json}: nitride run-specific ALIGNN configs.
  \item \texttt{results/oxide/N...\_seed.../dataset\_root/}: oxide dataset materialization for each run.
  \item \texttt{results/nitride/N...\_seed.../dataset\_root/}: nitride dataset materialization for each run.
  \item \texttt{results/.../finetune\_last2/summary.json}: primary run summary files.
  \item \texttt{results/.../finetune\_last2/history\_train.json}: training history for each run.
  \item \texttt{results/.../finetune\_last2/history\_val.json}: validation history for each run.
  \item \texttt{results/.../finetune\_last2/prediction\_results\_test\_set.csv}: per-structure test predictions.
  \item \texttt{reports/week2/finetune\_runs.csv}: one row per Week 2 run.
  \item \texttt{reports/week2/finetune\_summary\_by\_N.csv}: aggregate mean and standard deviation table.
  \item \texttt{reports/week2/finetune\_summary\_wide.csv}: wide-format comparison table across families and training sizes.
  \item \texttt{reports/week2/zero\_shot\_summary.csv}: zero-shot comparison table.
  \item \texttt{reports/week2/run\_suite\_summary.json}: summary of the additional seed batch execution.
\end{itemize}

\section{Verification and Reproducibility}
The Week 2 checker script is:

\begin{lstlisting}[language=bash]
scripts/shared/check_week2_status.sh .
\end{lstlisting}

This checker confirms:
\begin{itemize}
  \item all 36 expected \texttt{summary.json} files exist;
  \item all expected aggregate CSV and plot artifacts exist;
  \item every aggregate row reports exactly three seeds.
\end{itemize}

This makes the final Week 2 state independently auditable from the command line.

\section{Limitations and Next Steps}
The Week 2 study is complete, but the findings point naturally to the next project phase:
\begin{itemize}
  \item compare against from-scratch training to determine whether transfer is helping at all in this family-restricted regime;
  \item test whether deeper unfreezing improves transfer, especially for the nitride arm;
  \item continue with later-week embedding and representation analysis to study how oxide and nitride internal states differ after adaptation;
  \item if time permits, expand from three seeds to five seeds for tighter uncertainty estimates.
\end{itemize}

\section{Conclusion}
Week 2 is fully completed in this workspace at the level expected for a serious experimental record.
The repository now contains the completed oxide and nitride fine-tuning sweeps over all required training sizes, the corresponding three-seed summary tables, and the final learning-curve plots.

The key scientific message of Week 2 is clear:
oxide adaptation is smoother and more stable than nitride adaptation, but the current partial-fine-tuning setup does not yet outperform zero-shot for either family.
That result is still valuable because it narrows the likely bottlenecks and provides a strong motivation for the Week 3 from-scratch comparison.

\appendix
\section{Appendix: Per-Seed Week 2 Results}
The table below records the individual test MAE and best validation loss for every Week 2 run.
Including this appendix makes the report self-contained even without the CSV files.

\small
\begin{longtable}{lrrrrrr}
\toprule
Family & $N$ & Seed & Train & Val & Test MAE & Best val L1 \\
\midrule
\endhead
Oxide & 10 & 0 & 5 & 5 & 0.06102 & 0.04938 \\
Oxide & 10 & 1 & 5 & 5 & 0.04149 & 0.02279 \\
Oxide & 10 & 2 & 5 & 5 & 0.03604 & 0.00957 \\
Oxide & 50 & 0 & 45 & 5 & 0.07005 & 0.03863 \\
Oxide & 50 & 1 & 45 & 5 & 0.06449 & 0.05149 \\
Oxide & 50 & 2 & 45 & 5 & 0.05069 & 0.03481 \\
Oxide & 100 & 0 & 90 & 10 & 0.05258 & 0.02885 \\
Oxide & 100 & 1 & 90 & 10 & 0.05445 & 0.04228 \\
Oxide & 100 & 2 & 90 & 10 & 0.05114 & 0.05514 \\
Oxide & 200 & 0 & 180 & 20 & 0.05364 & 0.03136 \\
Oxide & 200 & 1 & 180 & 20 & 0.05373 & 0.03322 \\
Oxide & 200 & 2 & 180 & 20 & 0.04785 & 0.02908 \\
Oxide & 500 & 0 & 450 & 50 & 0.04711 & 0.02146 \\
Oxide & 500 & 1 & 450 & 50 & 0.04869 & 0.03870 \\
Oxide & 500 & 2 & 450 & 50 & 0.04665 & 0.02520 \\
Oxide & 1000 & 0 & 900 & 100 & 0.04637 & 0.02585 \\
Oxide & 1000 & 1 & 900 & 100 & 0.04565 & 0.03382 \\
Oxide & 1000 & 2 & 900 & 100 & 0.04466 & 0.02773 \\
Nitride & 10 & 0 & 5 & 5 & 0.11757 & 0.09397 \\
Nitride & 10 & 1 & 5 & 5 & 0.08263 & 0.03392 \\
Nitride & 10 & 2 & 5 & 5 & 0.09605 & 0.06398 \\
Nitride & 50 & 0 & 45 & 5 & 0.16176 & 0.15772 \\
Nitride & 50 & 1 & 45 & 5 & 0.16006 & 0.12946 \\
Nitride & 50 & 2 & 45 & 5 & 0.12267 & 0.10150 \\
Nitride & 100 & 0 & 90 & 10 & 0.25599 & 0.27827 \\
Nitride & 100 & 1 & 90 & 10 & 0.28827 & 0.18875 \\
Nitride & 100 & 2 & 90 & 10 & 0.16836 & 0.14364 \\
Nitride & 200 & 0 & 180 & 20 & 0.14084 & 0.14585 \\
Nitride & 200 & 1 & 180 & 20 & 0.25077 & 0.08844 \\
Nitride & 200 & 2 & 180 & 20 & 0.13412 & 0.08733 \\
Nitride & 500 & 0 & 450 & 50 & 0.11103 & 0.06797 \\
Nitride & 500 & 1 & 450 & 50 & 0.11399 & 0.06323 \\
Nitride & 500 & 2 & 450 & 50 & 0.10659 & 0.06509 \\
Nitride & 1000 & 0 & 900 & 100 & 0.09962 & 0.05783 \\
Nitride & 1000 & 1 & 900 & 100 & 0.09825 & 0.06330 \\
Nitride & 1000 & 2 & 900 & 100 & 0.10333 & 0.06613 \\
\bottomrule
\end{longtable}
\normalsize

\section{Appendix: Source Documents Used}
The Week 2 report was written from the executed workspace artifacts together with the project materials stored locally in the repository, especially:
\begin{itemize}
  \item the RES201 project brief in \texttt{Project\_Task/RES201 Project.pdf};
  \item the ALIGNN tutorial and project timeline notes in \texttt{Project\_Task/};
  \item the saved Week 2 result summaries in \texttt{reports/week2/};
  \item the per-run outputs in \texttt{results/oxide/} and \texttt{results/nitride/}.
\end{itemize}

\end{document}
